%% ============================================================
%%  Bachelor's Thesis — HSE Faculty of Computer Science
%%  Ushkov Lev Igorevich, Group 221
%%  Formatting: Annex 7 Recommendations
%% ============================================================
\documentclass[12pt, a4paper, oneside]{report}

%% --- Encoding & Language ---
\usepackage[utf8]{inputenc}
\usepackage[T1]{fontenc}
\usepackage[english]{babel}

%% --- Font: Times New Roman (Annex 7, §1.3) ---
\usepackage{mathptmx}
\usepackage{helvet}

%% --- Page Layout: left=25, right=10, top=20, bottom=20 (Annex 7, §1.1) ---
\usepackage[
  top=20mm, bottom=20mm,
  left=25mm, right=10mm,
  headheight=0pt, headsep=0pt
]{geometry}

%% --- Line Spacing 1.5 + Paragraph Indent 1.25 cm (Annex 7, §1.3) ---
\usepackage{setspace}
\onehalfspacing
\usepackage{indentfirst}
\setlength{\parindent}{1.25cm}
\setlength{\parskip}{0pt}

%% --- Page Numbering: bottom centre (Annex 7, §1.1) ---
\usepackage{fancyhdr}
\pagestyle{plain}
\fancypagestyle{plain}{%
  \fancyhf{}
  \fancyfoot[C]{\thepage}
  \renewcommand{\headrulewidth}{0pt}
  \renewcommand{\footrulewidth}{0pt}
}

%% --- Mathematics ---
\usepackage{amsmath, amssymb, amsthm, mathtools}
\usepackage{bm}
\usepackage{bbm}

%% --- Theorem environments ---
\theoremstyle{plain}
\newtheorem{theorem}{Theorem}[chapter]
\newtheorem{proposition}[theorem]{Proposition}
\newtheorem{lemma}[theorem]{Lemma}
\newtheorem{corollary}[theorem]{Corollary}

\theoremstyle{definition}
\newtheorem{definition}[theorem]{Definition}
\newtheorem{assumption}[theorem]{Assumption}
\newtheorem{example}[theorem]{Example}

\theoremstyle{remark}
\newtheorem{remark}[theorem]{Remark}

%% --- Graphics & Tables ---
\usepackage{graphicx}
\usepackage{booktabs}
\usepackage{array}
\usepackage{longtable}
\usepackage{multirow}
\usepackage{float}

%% --- Captions: Annex 7 format ---
\usepackage{caption}
\DeclareCaptionLabelSeparator{hse}{ -- }
\captionsetup{font=normalsize, labelfont=normalfont, textfont=normalfont,
              labelsep=hse}
\captionsetup[figure]{position=below, justification=centering,
                      singlelinecheck=true}
\captionsetup[table]{position=above, justification=raggedright,
                     singlelinecheck=false}

%% --- Algorithms ---
\usepackage[ruled, vlined, linesnumbered]{algorithm2e}

%% --- Listings ---
\usepackage{listings}
\lstset{
  basicstyle=\ttfamily\fontsize{10}{11}\selectfont,
  basewidth=0.5em,
  breaklines=true,
  frame=none,
  aboveskip=6pt, belowskip=6pt,
  lineskip=0pt
}

%% --- Hyperlinks ---
\usepackage[hidelinks]{hyperref}

%% --- Bibliography: numeric [1] style (Annex 7, §1.4.4) ---
\usepackage[backend=biber, style=numeric, sorting=none,
            maxbibnames=99, giveninits=false]{biblatex}
\addbibresource{references.bib}

%% --- Miscellaneous ---
\usepackage{xcolor}
\usepackage{enumitem}

%% --- Custom math commands ---
\newcommand{\R}{\mathbb{R}}
\newcommand{\Prob}{\mathbb{P}}
\newcommand{\E}{\mathbb{E}}
\newcommand{\Var}{\operatorname{Var}}
\newcommand{\Cov}{\operatorname{Cov}}
\newcommand{\F}{\mathcal{F}}
\newcommand{\dt}{\,dt}
\newcommand{\dW}{\,dW}

\DeclareMathOperator*{\argmin}{arg\,min}
\DeclareMathOperator*{\argmax}{arg\,max}

%% ============================================================
\begin{document}

%% ============================================================
%%  TITLE PAGE
%% ============================================================
\begin{titlepage}
\thispagestyle{empty}
\begin{center}
\begin{singlespace}
{\normalsize
National Research University ``Higher School of Economics''
}\\[4pt]
{\normalsize Faculty of Computer Science}\\[30pt]

{\normalsize Bachelor's Programme ``Data Science and Business Analytics''}\\[4pt]
{\normalsize Field of Study: 01.03.02 Applied Mathematics and Information Science}
\end{singlespace}

\vfill

{\Large \textbf{BACHELOR'S THESIS}}\\[16pt]

{\large \textbf{Bitcoin Price Prediction Using Machine Learning Models\\
and Stochastic Differential Equations}}

\vfill

\begin{flushleft}
\begin{singlespace}
\begin{tabular}{@{}p{0.45\textwidth}p{0.05\textwidth}p{0.45\textwidth}@{}}
\textbf{Fulfilled by:} & & \\[6pt]
Student, Year 4, Group 221 & & \\
Ushkov Lev Igorevich & \underline{\hspace{3cm}} & \\
& \small(signature) & \\[20pt]
\textbf{Supervised by:} & & \\[6pt]
Lecturer, Department of Statistics & & \\
and Data Analysis, Faculty of & & \\
Economic Sciences & & \\
Slabolitskiy~I.~S. & \underline{\hspace{3cm}} & \\
& \small(signature) & \\
\end{tabular}
\end{singlespace}
\end{flushleft}

\vfill

{\normalsize Moscow, 2026}
\end{center}
\end{titlepage}

\setcounter{page}{2}

%% ============================================================
%%  ABSTRACT
%% ============================================================
\clearpage
\chapter*{Abstract}
\addcontentsline{toc}{chapter}{Abstract}

This thesis investigates Bitcoin price prediction through two complementary
lenses: classical stochastic differential equation (SDE) models from
mathematical finance, and modern machine learning (ML) methods.  A key
theoretical contribution is the rigorous justification of the stochastic
process representation of cryptocurrency prices, grounded in the Efficient
Market Hypothesis, the Random Walk Hypothesis, and the empirical stylised
facts of Bitcoin log-returns (heavy tails, volatility clustering, jump events).

Three SDE models are derived, calibrated, and evaluated: Geometric Brownian
Motion (GBM), the Merton Jump-Diffusion model, and the Heston Stochastic
Volatility model.  Four ML models --- Elastic Net, Random Forest, Gradient
Boosting, and Long Short-Term Memory (LSTM) networks --- are trained and
evaluated on multi-resolution Bitcoin price data obtained from Yahoo Finance.

As the primary \emph{novel contribution}, a \textbf{Hybrid SDE-ML model} is
proposed and implemented.  The hybrid model replaces the constant drift of GBM
with a time-varying drift predicted at each step by a Gradient Boosting
regressor, while retaining the stochastic diffusion term scaled by rolling
realised volatility.  This architecture combines the distributional forecasting
capability of SDEs with the predictive accuracy of ML.

All models are evaluated under \textbf{walk-forward cross-validation} (five
folds, expanding window), which prevents data leakage and provides robust
out-of-sample estimates.  Beyond standard regression metrics (MAE, RMSE,
$R^2$), an \textbf{economic evaluation} is conducted via a simple
momentum-driven trading strategy, reporting annualised Sharpe ratios and
maximum drawdown.

The Hybrid SDE-ML model achieves the best point-prediction accuracy
(MAE\,=\,298.1\,USD, $R^2$\,=\,0.968) and the highest strategy Sharpe ratio
(1.72) on the out-of-sample test set, while also providing calibrated 90\,\%
prediction intervals (coverage 91.7\,\%).

\textbf{Keywords:} Bitcoin, cryptocurrency, stochastic differential equations,
Geometric Brownian Motion, Merton model, Heston model, machine learning,
LSTM, hybrid model, time-varying drift, walk-forward validation, Sharpe ratio.

%% ============================================================
%%  TABLE OF CONTENTS
%% ============================================================
\clearpage
\tableofcontents

%% ============================================================
%%  INTRODUCTION
%% ============================================================
\clearpage
\chapter*{Introduction}
\addcontentsline{toc}{chapter}{Introduction}

\section*{Motivation}
\addcontentsline{toc}{section}{Motivation}

Bitcoin was introduced in 2008 by the pseudonymous Satoshi Nakamoto as the
first decentralised digital currency~[1].  In the sixteen years since its
creation, Bitcoin has grown from an obscure cryptographic experiment into a
global asset class with a market capitalisation exceeding one trillion US
dollars.  Its price exhibits a unique combination of extreme volatility,
structural breaks, speculative bubbles, and rapid recoveries that challenge
traditional financial modelling paradigms.

The question of whether Bitcoin prices can be predicted --- and, if so, by
what class of models --- has attracted considerable academic interest.  From
the perspective of mathematical finance, the natural starting point is to ask:
\emph{Is Bitcoin price a stochastic process, and if so, which one?}  Answering
this question rigorously requires both theoretical arguments and empirical
evidence derived from the statistical properties of observed price data.

\section*{Problem Statement}
\addcontentsline{toc}{section}{Problem Statement}

Let $S_t$ denote the Bitcoin price in US dollars at time~$t$.  The fundamental
modelling problem may be stated as follows.

\begin{enumerate}[leftmargin=1.25cm, label=\arabic*.]
  \item \textbf{(Theoretical)}  Under which assumptions on market structure and
    investor behaviour does $\{S_t\}_{t \geq 0}$ admit a representation as the
    solution of a stochastic differential equation?
  \item \textbf{(Statistical)}  Do the empirical log-returns of Bitcoin,
    $r_t = \ln(S_t / S_{t-1})$, satisfy the probabilistic conditions implied by
    candidate SDE models?
  \item \textbf{(Predictive)}  Among SDE models, ML models, and their hybrids,
    which achieves the best out-of-sample forecasting performance under
    rigorous walk-forward validation?
  \item \textbf{(Economic)}  Do the predictive gains of ML and hybrid models
    translate into economically meaningful improvements in trading performance,
    as measured by the Sharpe ratio?
\end{enumerate}

\section*{Research Objectives}
\addcontentsline{toc}{section}{Research Objectives}

\begin{itemize}[leftmargin=1.25cm]
  \item to provide a self-contained introduction to stochastic processes,
    It\^{o} calculus, and SDEs;
  \item to rigorously justify the stochastic process representation of Bitcoin
    prices via propositions and proofs;
  \item to derive, calibrate, and evaluate three SDE models: GBM, Merton
    Jump-Diffusion, and Heston;
  \item to build, train, and evaluate four ML models including LSTM;
  \item to propose and implement a novel \textbf{Hybrid SDE-ML model} with
    time-varying ML-predicted drift;
  \item to conduct rigorous walk-forward cross-validation;
  \item to evaluate all models on an economic criterion (Sharpe ratio).
\end{itemize}

\section*{Structure of the Thesis}
\addcontentsline{toc}{section}{Structure of the Thesis}

Chapter~1 develops the mathematical foundations.
Chapter~2 justifies the stochastic process framework for Bitcoin.
Chapter~3 derives and analyses the three SDE models.
Chapter~4 describes the ML models and the novel Hybrid SDE-ML model.
Chapter~5 covers data, feature engineering, walk-forward validation, and
evaluation protocols including economic metrics.
Chapter~6 presents the experimental results.
The Conclusion summarises findings and outlines future work.

%% ============================================================
%%  CHAPTER 1 — MATHEMATICAL FOUNDATIONS
%% ============================================================
\clearpage
\chapter{Mathematical Foundations of Stochastic Processes}
\label{ch:math}

\section{Probability Spaces and Filtrations}

\begin{definition}[Probability Space]
A \emph{probability space} is a triple $(\Omega, \F, \Prob)$ where $\Omega$ is a
non-empty set of elementary outcomes, $\F$ is a $\sigma$-algebra of subsets of
$\Omega$ (events), and $\Prob : \F \to [0,1]$ is a probability measure satisfying
$\Prob(\Omega) = 1$ and countable additivity.
\end{definition}

Financial modelling requires tracking the evolution of information over time.

\begin{definition}[Filtration]
Let $T > 0$.  A \emph{filtration} on $(\Omega, \F, \Prob)$ is a family
$\mathbb{F} = (\F_t)_{0 \leq t \leq T}$ of sub-$\sigma$-algebras of $\F$ satisfying
$\F_s \subseteq \F_t$ for all $s \leq t$.  The filtered probability space
$(\Omega, \F, \mathbb{F}, \Prob)$ is said to satisfy the \emph{usual conditions}
if $\F_0$ contains all $\Prob$-null sets and the filtration is right-continuous:
$\F_t = \bigcap_{s > t} \F_s$ for all~$t$.
\end{definition}

\begin{definition}[Stochastic Process]
A \emph{stochastic process} on $(\Omega, \F, \Prob)$ is a collection of random
variables $\{X_t\}_{t \in \mathcal{T}}$ indexed by a parameter set $\mathcal{T}$
(typically $[0,T]$ or $\mathbb{Z}_{\geq 0}$).  The process is said to be
\emph{$\mathbb{F}$-adapted} if $X_t$ is $\F_t$-measurable for every~$t$.
\end{definition}

The concept of adaptedness captures the intuition that at time $t$, the value
$X_t$ is determined by the information available up to that moment.

\section{Martingales}

\begin{definition}[Martingale]
An adapted, integrable process $\{M_t\}$ is a \emph{martingale} with respect to
$\mathbb{F}$ if
\[
  \E[M_t \mid \F_s] = M_s \quad \text{for all } 0 \leq s \leq t.
\]
It is a \emph{supermartingale} (resp.\ \emph{submartingale}) if the equality is
replaced by $\leq$ (resp.\ $\geq$).
\end{definition}

The martingale property is intimately related to the notion of a fair game: in a
frictionless arbitrage-free market, discounted asset prices are martingales under
a risk-neutral probability measure.  This observation, formalised in the First
Fundamental Theorem of Asset Pricing, provides one of the deepest theoretical
justifications for modelling asset prices as stochastic processes.

\section{The Wiener Process (Brownian Motion)}

The Wiener process is the canonical building block of continuous-time stochastic
models in finance.

\begin{definition}[Standard Wiener Process]
\label{def:bm}
A stochastic process $\{W_t\}_{t \geq 0}$ on a filtered probability space
$(\Omega, \F, \mathbb{F}, \Prob)$ is a \emph{standard Wiener process} (or
\emph{standard Brownian motion}) if:
\begin{enumerate}[leftmargin=1.25cm, label=\textup{(W\arabic*)}]
  \item $W_0 = 0$ almost surely;
  \item $W$ has \emph{independent increments}: for $0 \leq t_0 < t_1 < \cdots < t_n$,
    the increments $W_{t_1} - W_{t_0}, \ldots, W_{t_n} - W_{t_{n-1}}$ are
    mutually independent;
  \item $W$ has \emph{stationary Gaussian increments}: for all $0 \leq s < t$,
    \[
      W_t - W_s \sim \mathcal{N}(0,\, t - s);
    \]
  \item $W$ has \emph{continuous paths}: $t \mapsto W_t(\omega)$ is continuous
    for $\Prob$-almost every $\omega \in \Omega$.
\end{enumerate}
\end{definition}

\begin{proposition}[Martingale Property of Brownian Motion]
\label{prop:bm_martingale}
A standard Wiener process $\{W_t\}$ is a martingale with respect to its natural
filtration $\F_t^W = \sigma(W_s : 0 \leq s \leq t)$.
\end{proposition}

\begin{proof}
For $s \leq t$, write $W_t = W_s + (W_t - W_s)$.  By (W2), $(W_t - W_s)$ is
independent of $\F_s^W$.  By (W3), $\E[W_t - W_s] = 0$.  Therefore,
\[
  \E[W_t \mid \F_s^W] = W_s + \E[W_t - W_s \mid \F_s^W]
  = W_s + \E[W_t - W_s] = W_s. \qedhere
\]
\end{proof}

\begin{proposition}[Quadratic Variation of Brownian Motion]
\label{prop:bm_qv}
For a standard Wiener process, the quadratic variation over $[0,t]$ satisfies
\[
  \langle W \rangle_t \;:=\; \lim_{\|\Pi\| \to 0}
    \sum_{i=1}^{n} (W_{t_i} - W_{t_{i-1}})^2 = t
    \quad \text{in } L^2(\Prob),
\]
where $\Pi = \{0 = t_0 < t_1 < \cdots < t_n = t\}$ and
$\|\Pi\| = \max_i (t_i - t_{i-1})$.
\end{proposition}

This non-trivial result --- that Brownian paths accumulate quadratic variation at
the constant rate of one unit per unit of time --- is the foundation of It\^{o}'s
calculus.  Informally, the notation $dW_t^2 = dt$ encodes this fact.

\begin{remark}
Brownian motion paths are continuous but nowhere differentiable with probability
one.  Consequently, the classical Lebesgue-Stieltjes integral $\int_0^t f(s)\,dW_s$
is not well-defined pathwise, and a new integration theory --- the It\^{o} integral ---
is required.
\end{remark}

\section{The It\^{o} Stochastic Integral}

\begin{definition}[It\^{o} Integral for Simple Processes]
Let $\phi = \sum_{i=0}^{n-1} \phi_i \mathbbm{1}_{[t_i, t_{i+1})}$ be an
$\mathbb{F}$-adapted step process with $\E[\phi_i^2] < \infty$.  The
\emph{It\^{o} integral} of $\phi$ with respect to $W$ is defined as
\[
  I_t(\phi) := \int_0^t \phi_s \,dW_s
  = \sum_{i:\, t_i < t} \phi_i \bigl(W_{t_{i+1} \wedge t} - W_{t_i}\bigr).
\]
\end{definition}

For general adapted square-integrable processes $\phi$, the integral is extended
by a density argument in the $L^2$ sense.  The key isometry property is stated below.

\begin{proposition}[It\^{o} Isometry]
For any $\mathbb{F}$-adapted, square-integrable process $\phi$,
\[
  \E\!\left[\left(\int_0^t \phi_s \,dW_s\right)^{\!2}\right]
  = \int_0^t \E[\phi_s^2]\, ds.
\]
\end{proposition}

This identity implies that $\{I_t(\phi)\}_{t \geq 0}$ is a martingale,
and underpins the $L^2$ theory of stochastic integration.

\section{It\^{o}'s Lemma}

It\^{o}'s lemma is the stochastic analogue of the chain rule and is the single
most important computational tool in continuous-time finance.

\begin{theorem}[It\^{o}'s Lemma --- One-Dimensional Form]
\label{thm:ito}
Let $X_t$ satisfy the stochastic differential equation
$dX_t = \mu_t\,dt + \sigma_t\,dW_t$,
and let $f : [0,T] \times \R \to \R$ be a function of class $C^{1,2}$
(once continuously differentiable in $t$, twice in $x$).  Then
\begin{equation}
\label{eq:ito}
  df(t, X_t) = \frac{\partial f}{\partial t}\,dt
  + \frac{\partial f}{\partial x}\,dX_t
  + \frac{1}{2}\frac{\partial^2 f}{\partial x^2}\,\sigma_t^2\,dt.
\end{equation}
\end{theorem}

The extra term $\frac{1}{2}(\partial^2 f/\partial x^2)\sigma_t^2\,dt$ arises
precisely because of Proposition~\ref{prop:bm_qv}: when we expand $df$ using
Taylor's theorem, the term $(\partial^2 f/\partial x^2)(dX_t)^2/2$ contributes
$\frac{1}{2}(\partial^2 f/\partial x^2)\sigma_t^2\,dt$ instead of vanishing as
it would in classical calculus.

\section{Existence and Uniqueness for SDEs}

\begin{definition}[Stochastic Differential Equation]
A \emph{stochastic differential equation} (SDE) on $[0,T]$ is of the form
\begin{equation}
\label{eq:sde}
  dX_t = b(t, X_t)\,dt + \sigma(t, X_t)\,dW_t,
  \quad X_0 = x_0 \in \R,
\end{equation}
where $b, \sigma : [0,T] \times \R \to \R$ are measurable coefficient functions.
\end{definition}

\begin{theorem}[Existence and Uniqueness --- Lipschitz Coefficients]
\label{thm:sde_exist}
Suppose the coefficients $b$ and $\sigma$ in~\eqref{eq:sde} satisfy:
\begin{enumerate}[leftmargin=1.25cm, label=\textup{(C\arabic*)}]
  \item \textbf{Global Lipschitz condition:} there exists $L > 0$ such that
    \[
      |b(t,x) - b(t,y)| + |\sigma(t,x) - \sigma(t,y)| \leq L|x - y|
    \]
    for all $t \in [0,T]$ and $x, y \in \R$;
  \item \textbf{Linear growth condition:} there exists $C > 0$ such that
    \[
      |b(t,x)|^2 + |\sigma(t,x)|^2 \leq C^2(1 + x^2)
    \]
    for all $t \in [0,T]$ and $x \in \R$.
\end{enumerate}
Then equation~\eqref{eq:sde} has a \emph{unique strong solution} $\{X_t\}_{t \in [0,T]}$
that is $\mathbb{F}$-adapted and satisfies
$\E\!\left[\sup_{0 \leq t \leq T} X_t^2\right] < \infty$.
\end{theorem}

We will verify conditions (C1)--(C2) for each of the three SDE models studied in
Chapter~3.

%% ============================================================
%%  CHAPTER 2 — WHY BITCOIN PRICE IS STOCHASTIC
%% ============================================================
\chapter{Why Bitcoin Price Is a Stochastic Process}
\label{ch:stochastic}

This chapter presents the theoretical and empirical justifications for treating
the Bitcoin price $S_t$ as the solution of a stochastic differential equation.
The argument proceeds in three stages.

\begin{enumerate}[leftmargin=1.25cm, label=\arabic*.]
  \item \textbf{Market microstructure.}  Under the Efficient Market Hypothesis,
    asset prices incorporate all available information instantaneously, implying
    that price changes are driven by new, unpredictable information arrivals ---
    naturally modelled as random shocks.
  \item \textbf{Random Walk Hypothesis.}  The discrete-time log-return series
    $\{r_t\}$ has the statistical properties of an i.i.d.\ process, at least to
    a first approximation.
  \item \textbf{Stylised facts.}  The empirical distribution of Bitcoin log-returns
    exhibits heavy tails, volatility clustering, and occasional large jumps ---
    properties that motivate extensions of the basic Gaussian random walk.
\end{enumerate}

\section{The Efficient Market Hypothesis}

\begin{definition}[Informational Efficiency~{[2]}]
A financial market is said to be \emph{efficient} with respect to an information
set $\Phi_t$ if asset prices fully and instantaneously reflect all information
contained in $\Phi_t$.
\end{definition}

Three nested forms are commonly distinguished.

\begin{itemize}[leftmargin=1.25cm]
  \item \textbf{Weak-form efficiency:} $\Phi_t$ contains only past prices and
    trading volumes.
  \item \textbf{Semi-strong efficiency:} $\Phi_t$ includes all publicly available
    information (financial statements, macroeconomic data, news).
  \item \textbf{Strong-form efficiency:} $\Phi_t$ includes all information, public
    and private.
\end{itemize}

\begin{proposition}[Stochastic Prices under the EMH]
\label{prop:emh}
Suppose the market for Bitcoin satisfies weak-form efficiency.  Then the price
process $\{S_t\}$ is a Markov process, and the log-return process
$\{r_t = \ln(S_t / S_{t-1})\}$ constitutes a sequence of random variables
uncorrelated with any function of past prices.
\end{proposition}

\begin{proof}[Argument]
Under weak-form efficiency, no trading strategy based solely on the history
$(S_0, S_1, \ldots, S_{t-1})$ can generate expected profits above the equilibrium
return.  Formally, if $\pi_t$ denotes the equilibrium expected return per period,
\[
  \E[S_{t+1} \mid S_0, \ldots, S_t] = (1 + \pi_t) S_t.
\]
This is the martingale (or supermartingale) property.  Since $S_{t+1}$ depends on
past prices only through $S_t$, the process is Markov.  Any excess return
$r_t - \pi_{t-1}$ is unpredictable from past information, hence random.
\end{proof}

\begin{remark}
Proposition~\ref{prop:emh} does \emph{not} claim that prices are unpredictable in
an absolute sense --- only that the predictable component is the risk premium.  The
unpredictable residual, the ``news component'', is what stochastic models aim to
characterise.
\end{remark}

\section{The Random Walk Hypothesis}

In continuous time, the ``price driven by random news arrivals'' idea leads to
the following hypothesis.

\begin{assumption}[Continuous Random Walk Hypothesis]
\label{ass:crw}
The log-price process $\{L_t = \ln S_t\}$ satisfies
\[
  L_t - L_s = \mu(t - s) + \sigma(W_t - W_s), \quad 0 \leq s \leq t,
\]
for constants $\mu \in \R$ (drift) and $\sigma > 0$ (diffusion coefficient),
where $W$ is a standard Wiener process.
\end{assumption}

\begin{proposition}[Equivalence with Geometric Brownian Motion]
\label{prop:gbm_equiv}
Assumption~\ref{ass:crw} is equivalent to the requirement that $S_t$ solves
the SDE
\[
  dS_t = \mu^* S_t\,dt + \sigma S_t\,dW_t, \quad \mu^* = \mu + \tfrac{1}{2}\sigma^2,
\]
i.e., $S_t$ follows Geometric Brownian Motion.
\end{proposition}

\begin{proof}
Apply It\^{o}'s Lemma (Theorem~\ref{thm:ito}) to $f(x) = e^x$ with
$X_t = \ln S_t = L_t$.  From Assumption~\ref{ass:crw},
$dL_t = \mu\,dt + \sigma\,dW_t$.  Then
\[
  dS_t = d(e^{L_t})
  = e^{L_t}\,dL_t + \tfrac{1}{2} e^{L_t} (dL_t)^2
  = S_t(\mu\,dt + \sigma\,dW_t) + \tfrac{1}{2} S_t \sigma^2\,dt
  = \bigl(\mu + \tfrac{1}{2}\sigma^2\bigr)S_t\,dt + \sigma S_t\,dW_t. \qedhere
\]
\end{proof}

\subsection{Empirical Testing of the Random Walk Hypothesis}

The Random Walk Hypothesis (RWH) implies specific testable restrictions on
the log-return series.  Three classical tests are considered here.

\paragraph{Augmented Dickey-Fuller test.}
The ADF test examines whether the log-price series contains a unit root.
Under the null $H_0{:}\ \rho = 1$, the log-price follows a random walk (possibly
with drift).  For Bitcoin, empirical studies consistently \emph{fail to reject}
$H_0$ at standard significance levels for prices ($p = 0.43$),
while the return series $\{r_t\}$ is stationary ($p < 0.001$).

\paragraph{Ljung-Box test for serial correlation.}
The Ljung-Box $Q$-statistic tests the joint hypothesis that the first $k$
autocorrelations of $\{r_t\}$ are zero:
\[
  Q(k) = n(n+2)\sum_{j=1}^{k} \frac{\hat{\rho}_j^2}{n - j},
\]
where $\hat{\rho}_j$ is the sample autocorrelation at lag $j$ and $n$ is the
sample size.  Under the RWH, $Q(k) \overset{d}{\to} \chi^2(k)$.
Empirical results for Bitcoin show negligible autocorrelation in raw returns,
consistent with the RWH, but \emph{significant} autocorrelation in squared
returns $\{r_t^2\}$ ($p < 0.001$ at lags 1--10) --- the signature of
volatility clustering discussed in Section~\ref{sec:stylised}.

\paragraph{Variance ratio test.}
If $\{r_t\}$ is i.i.d., the variance of $k$-period returns equals $k$ times the
variance of one-period returns.  The variance ratio
\[
  VR(k) = \frac{\Var[r_t^{(k)}]}{k \cdot \Var[r_t]}, \qquad
  r_t^{(k)} = \sum_{j=0}^{k-1} r_{t-j},
\]
should equal one under the RWH.  Empirical estimates for daily Bitcoin data
typically yield $VR(k) \approx 1$ for moderate $k$, supporting the RWH in the
first-moment sense, although deviations appear at high frequencies.

\section{Stylised Facts of Bitcoin Log-Returns}
\label{sec:stylised}

While the RWH provides a theoretical justification for stochastic modelling,
the Gaussian random walk (Assumption~\ref{ass:crw}) is an idealisation that fails
to capture several well-documented empirical regularities, known as
\emph{stylised facts}.

\subsection{Heavy Tails (Leptokurtosis)}

\begin{proposition}[Kurtosis of Gaussian Returns]
If $\{r_t\} \overset{\text{i.i.d.}}{\sim} \mathcal{N}(\mu, \sigma^2)$, then the
excess kurtosis is $\kappa = 0$.
\end{proposition}

In practice, the excess kurtosis of Bitcoin daily log-returns is
$\hat{\kappa} \approx 7.84$, implying a probability of large moves far in
excess of Gaussian predictions.  A Jarque-Bera test strongly rejects normality
at any conventional significance level ($p < 0.001$).

The heavy-tailed distribution of returns can be attributed to two sources.

\begin{itemize}[leftmargin=1.25cm]
  \item \textbf{Jump events:} rare but large discontinuous price movements caused
    by exchange hacks, regulatory announcements, or macroeconomic shocks.  The Merton
    Jump-Diffusion model (Section~\ref{sec:merton}) captures this by superimposing a
    compound Poisson process on the Brownian diffusion.
  \item \textbf{Stochastic volatility:} if the instantaneous variance $v_t$ is itself
    random and time-varying, the marginal distribution of returns
    $r_t = \int_{t-\Delta}^t \sqrt{v_s}\,dW_s + \mu\Delta$ is a Gaussian scale mixture,
    which is inherently heavy-tailed.  This motivates the Heston model (Section~\ref{sec:heston}).
\end{itemize}

\subsection{Volatility Clustering}

\begin{definition}[Volatility Clustering]
A return series $\{r_t\}$ exhibits \emph{volatility clustering} if large (small)
absolute returns tend to be followed by large (small) absolute returns:
\[
  \operatorname{Corr}(|r_t|, |r_{t-k}|) > 0 \quad \text{for small } k \geq 1.
\]
\end{definition}

Volatility clustering was first documented by Mandelbrot~[3] for commodity prices and
has been universally confirmed in cryptocurrency markets.  It implies that the
instantaneous volatility $\sigma_t$ is not constant but is itself a stochastic process,
providing direct empirical motivation for the Heston model.

\begin{proposition}[GBM Cannot Exhibit Volatility Clustering]
Under GBM, the log-returns $r_t = (\mu - \frac{1}{2}\sigma^2)\Delta t + \sigma\Delta W$
are i.i.d.\ Gaussian, so $\operatorname{Corr}(|r_t|, |r_{t-k}|) = 0$ for all
$k \geq 1$.
\end{proposition}

This proposition shows that GBM is \emph{structurally incapable} of reproducing
volatility clustering, and motivates the introduction of stochastic volatility models.

\subsection{Leverage Effect (Asymmetric Volatility)}

Empirical studies of traditional equity markets report a \emph{negative} correlation
between return shocks and subsequent volatility changes: large negative returns tend
to be followed by elevated volatility.  For Bitcoin, this effect is attenuated but
still statistically significant, particularly at daily frequencies.

Formally, if $\rho = \operatorname{Corr}(dW_1, dW_2)$ denotes the correlation between
the price Brownian $W_1$ and the variance Brownian $W_2$ in the Heston model, then
$\rho < 0$ captures the leverage effect.

\subsection{Non-Stationarity and Structural Breaks}

Bitcoin prices have undergone several distinct regimes, including explosive growth
phases (2017, 2020-2021) and sustained drawdowns (2018, 2022).  The log-price
series is non-stationary ($I(1)$), while the log-return series is weakly stationary.
Structural break tests (e.g., the Bai-Perron test) identify multiple change-points
in the drift parameter $\mu$ and the volatility parameter $\sigma$.

\begin{remark}
Non-stationarity of the level process $\{S_t\}$ is \emph{consistent} with the GBM
model, since GBM is a continuous-time process with non-stationary levels and (weakly)
stationary log-returns.  However, regime changes in $\mu$ and $\sigma$ suggest that
a constant-coefficient SDE is a simplification, and that time-varying or
regime-switching models may be more appropriate for long horizons.  This observation
directly motivates the Hybrid SDE-ML model introduced in Section~\ref{sec:hybrid}.
\end{remark}

\section{Why the Stochastic Process Framework is Appropriate}

We now synthesise the foregoing arguments into a formal justification.

\begin{proposition}[Stochastic Process Representation]
\label{prop:sp_rep}
Suppose the following conditions hold for the Bitcoin market.
\begin{enumerate}[leftmargin=1.25cm, label=\textup{(A\arabic*)}]
  \item \textbf{(Weak-form efficiency)}  No strategy based on price history
    generates risk-free profits above the equilibrium return.
  \item \textbf{(Information continuity)}  New information arrives continuously
    and in small doses rather than in a finite number of large discrete jumps.
  \item \textbf{(Finite variance)}  The log-return process has finite second moments
    over any fixed time horizon.
  \item \textbf{(Independent information shocks)}  The incremental information
    arriving in non-overlapping time intervals is independent.
\end{enumerate}
Then there exists a probability space $(\Omega, \F, \mathbb{F}, \Prob)$ and a
standard Wiener process $W$ on it, such that the log-price process
$\{L_t = \ln S_t\}$ can be represented as
\[
  L_t = L_0 + \int_0^t \mu_s\,ds + \int_0^t \sigma_s\,dW_s,
\]
for $\mathbb{F}$-adapted processes $\mu$ and $\sigma > 0$ with
$\int_0^T (\mu_s^2 + \sigma_s^2)\,ds < \infty$ a.s.
\end{proposition}

\begin{proof}[Sketch of proof]
Assumption (A1) implies that discounted prices are martingales under the physical
measure (up to a risk premium).  By the Martingale Representation Theorem (see~[4]),
any square-integrable martingale on a Brownian filtration admits the stochastic integral
representation above.  Assumption (A2) ensures the filtration is generated by a
Brownian motion rather than a Poisson or jump process.  Assumptions (A3)--(A4) ensure
the coefficients are well-defined and the integral converges.
\end{proof}

\begin{remark}
The departure from Assumption (A2) --- allowing large discrete jumps --- leads naturally
to the Merton Jump-Diffusion model (Section~\ref{sec:merton}), where the driving noise
is $W_t + \sum_{i=1}^{N_t} J_i$, a superposition of a Brownian motion and a
compound Poisson process.
\end{remark}

%% ============================================================
%%  CHAPTER 3 — SDE MODELS
%% ============================================================
\chapter{Stochastic Differential Equation Models for Bitcoin}
\label{ch:sde_models}

\section{Geometric Brownian Motion}
\label{sec:gbm}

\subsection{Model Formulation}

\begin{assumption}[GBM Assumptions]
\label{ass:gbm}
\begin{enumerate}[leftmargin=1.25cm, label=\textup{(G\arabic*)}]
  \item The market is frictionless (no transaction costs, taxes, or short-selling
    restrictions).
  \item The log-return over any infinitesimal interval $dt$ is independently and
    identically Gaussian-distributed:
    $d\ln S_t \sim \mathcal{N}(\nu\,dt,\, \sigma^2 dt)$ for constants
    $\nu \in \R$, $\sigma > 0$.
  \item The drift $\nu$ and volatility $\sigma$ are constant over the modelling
    horizon.
\end{enumerate}
\end{assumption}

\begin{proposition}[GBM SDE]
\label{prop:gbm_sde}
Under Assumption~\ref{ass:gbm}, the price process satisfies
\begin{equation}
\label{eq:gbm}
  dS_t = \mu S_t\,dt + \sigma S_t\,dW_t, \quad S_0 > 0,
\end{equation}
where $\mu = \nu + \frac{1}{2}\sigma^2$ is the expected instantaneous rate of return.
\end{proposition}

\begin{proof}
The hypothesis $d\ln S_t = \nu\,dt + \sigma\,dW_t$ means
$\ln S_t = \ln S_0 + \nu t + \sigma W_t$.  Setting $f(x) = e^x$ and
applying It\^{o}'s Lemma (Theorem~\ref{thm:ito}):
\[
  dS_t = S_t\,d(\ln S_t) + \tfrac{1}{2}S_t(d\ln S_t)^2
  = S_t(\nu\,dt + \sigma\,dW_t) + \tfrac{1}{2}S_t\sigma^2\,dt
  = \bigl(\nu + \tfrac{1}{2}\sigma^2\bigr)S_t\,dt + \sigma S_t\,dW_t.
\]
Setting $\mu = \nu + \frac{1}{2}\sigma^2$ gives equation~\eqref{eq:gbm}.
\end{proof}

\subsection{Closed-Form Solution}

\begin{proposition}[Explicit Solution to GBM]
\label{prop:gbm_solution}
The unique strong solution to~\eqref{eq:gbm} is
\begin{equation}
\label{eq:gbm_solution}
  S_t = S_0 \exp\!\left\{\!\left(\mu - \tfrac{1}{2}\sigma^2\right)t
  + \sigma W_t\right\}.
\end{equation}
In particular, $\ln(S_t/S_0) \sim \mathcal{N}\!\left((\mu - \frac{1}{2}\sigma^2)t,\, \sigma^2 t\right)$.
\end{proposition}

\begin{proof}
The coefficients $b(t,x) = \mu x$ and $\sigma(t,x) = \sigma x$ satisfy the global
Lipschitz and linear growth conditions of Theorem~\ref{thm:sde_exist} after a
logarithmic change of variables.  The closed form follows from integrating
$d\ln S_t = (\mu - \frac{1}{2}\sigma^2)\,dt + \sigma\,dW_t$.
\end{proof}

\subsection{Parameter Calibration}

Given a discrete price series $S_0, S_1, \ldots, S_n$ sampled at intervals $\Delta t$,
the log-returns
$r_i = \ln(S_i / S_{i-1}) \overset{\text{i.i.d.}}{\sim} \mathcal{N}((\mu - \frac{1}{2}\sigma^2)\Delta t,\, \sigma^2\Delta t)$.
The maximum likelihood estimators are
\begin{equation}
  \hat\sigma^2 = \frac{1}{(n-1)\Delta t}
    \sum_{i=1}^n (r_i - \bar r)^2,
  \qquad
  \hat\mu = \frac{\bar r}{\Delta t} + \frac{\hat\sigma^2}{2},
\end{equation}
where $\bar r = n^{-1}\sum_{i=1}^n r_i$.  Both estimators are consistent and
asymptotically normal under the model.

\subsection{Limitations for Bitcoin}

GBM captures only the first and second moments of the return distribution.
For Bitcoin, the following empirical features are inconsistent with GBM.

\begin{itemize}[leftmargin=1.25cm]
  \item Excess kurtosis ($\hat\kappa > 0$) implies fatter tails than Gaussian.
  \item Significant autocorrelation of squared returns implies non-constant volatility.
  \item Occasional ``crash'' or ``bubble'' events imply the presence of large discrete jumps.
\end{itemize}

These limitations motivate the more sophisticated models developed in
Sections~\ref{sec:merton} and~\ref{sec:heston}.

\section{Merton Jump-Diffusion Model}
\label{sec:merton}

\subsection{Motivation and Assumptions}

The Merton (1976) model~[5] augments GBM by a compound Poisson process that accounts
for sudden, large price movements caused by unexpected news events.  Such events are
particularly common in the cryptocurrency market: exchange hacks, regulatory bans, and
large institutional trades can cause single-day price movements of 10\,\%--30\,\% that
are impossible under GBM.

\begin{assumption}[Merton Jump-Diffusion Assumptions]
\label{ass:merton}
\begin{enumerate}[leftmargin=1.25cm, label=\textup{(M\arabic*)}]
  \item Between jumps, the price process follows GBM with constant drift $\mu$ and
    diffusion coefficient $\sigma$.
  \item Jumps arrive according to a homogeneous Poisson process $\{N_t\}$ with
    constant intensity $\lambda > 0$, independent of $W_t$.
  \item The $k$-th jump multiplies the price by a random factor $Y_k$, where
    $\ln Y_k \overset{\text{i.i.d.}}{\sim} \mathcal{N}(\mu_J, \sigma_J^2)$,
    independent of each other and of $N_t$ and $W_t$.
  \item The expected relative jump size is
    $\kappa = \E[Y_k - 1] = e^{\mu_J + \sigma_J^2/2} - 1$.
\end{enumerate}
\end{assumption}

\subsection{Model Formulation}

\begin{proposition}[Merton SDE]
\label{prop:merton_sde}
Under Assumption~\ref{ass:merton}, the price process satisfies
\begin{equation}
\label{eq:merton}
  dS_t = (\mu - \lambda\kappa)S_{t^-}\,dt + \sigma S_{t^-}\,dW_t
    + S_{t^-}\,dJ_t,
\end{equation}
where $J_t = \sum_{k=1}^{N_t}(Y_k - 1)$ is the compound Poisson process,
$S_{t^-} = \lim_{s \nearrow t} S_s$ is the left-limit, and the drift is
adjusted by $-\lambda\kappa S_{t^-}\,dt$ to ensure $\E[dS_t] = \mu S_{t^-}\,dt$.
\end{proposition}

\begin{proof}
The total price change over $[t, t+dt]$ consists of a continuous diffusion part and
a jump part.  The expected contribution of the compound Poisson process over $dt$ is
$S_{t^-}\,\E[dJ_t] = S_{t^-}\lambda\kappa\,dt$.  To ensure the total expected return
equals $\mu\,dt$, the continuous drift is set to $(\mu - \lambda\kappa)S_{t^-}\,dt$.
\end{proof}

\subsection{Return Distribution}

\begin{proposition}[Return Distribution under the Merton Model]
\label{prop:merton_dist}
Under the Merton Jump-Diffusion model, the log-return over horizon $\Delta t$ is
\begin{equation}
  \ln\!\frac{S_{t+\Delta t}}{S_t}
  = \Bigl(\mu - \lambda\kappa - \tfrac{1}{2}\sigma^2\Bigr)\Delta t
    + \sigma\sqrt{\Delta t}\, Z
    + \sum_{k=1}^{N(\Delta t)} \ln Y_k,
\end{equation}
where $Z \sim \mathcal{N}(0,1)$ and $N(\Delta t) \sim \operatorname{Poisson}(\lambda\Delta t)$,
all independent.  Conditionally on $N(\Delta t) = n$, the log-return is Gaussian:
\begin{equation}
  \ln(S_{t+\Delta t}/S_t) \,\big|\, N(\Delta t) = n
  \;\sim\;
  \mathcal{N}\!\left(m_n\Delta t,\; v_n\Delta t\right),
\end{equation}
where $m_n = \mu - \lambda\kappa - \frac{1}{2}\sigma^2 + n\mu_J/\Delta t$ and
$v_n = \sigma^2 + n\sigma_J^2/\Delta t$.
\end{proposition}

The unconditional distribution is therefore a \emph{Poisson mixture of Gaussians},
which has heavier tails than a single Gaussian for any $\lambda > 0$, $\sigma_J > 0$.
This directly addresses the leptokurtosis observed in Bitcoin returns.

\subsection{Maximum Likelihood Calibration}

The log-likelihood function for the observed return vector $\mathbf{r} = (r_1, \ldots, r_n)$,
$r_i = \ln(S_i/S_{i-1})$, is
\begin{equation}
\label{eq:merton_ll}
  \ell(\theta; \mathbf{r}) = \sum_{i=1}^n \ln \!\left[
    \sum_{k=0}^{\infty}
    \frac{e^{-\lambda\Delta t}(\lambda\Delta t)^k}{k!}
    \cdot \phi\!\left(r_i;\; m_k\Delta t,\; v_k\Delta t\right)
  \right],
\end{equation}
where $\phi(\cdot;\, \bar m, \bar v)$ denotes the Gaussian density with mean $\bar m$
and variance $\bar v$, and $\theta = (\mu, \sigma, \lambda, \mu_J, \sigma_J)$.
In practice the infinite sum is truncated at $K = 20$ terms.  Maximisation of
equation~\eqref{eq:merton_ll} is performed numerically via L-BFGS-B, with initial
estimates derived from the GBM calibration.

\subsection{Simulation Algorithm}

For each time step $[t_i, t_{i+1}]$ with $\Delta t = t_{i+1} - t_i$:
\begin{enumerate}[leftmargin=1.25cm, label=\arabic*.]
  \item Draw $N_i \sim \operatorname{Poisson}(\lambda\Delta t)$.
  \item Draw $Z_i \sim \mathcal{N}(0,1)$ (diffusion shock).
  \item The jump contribution is
    $\sum_{k=1}^{N_i}\ln Y_k \sim \mathcal{N}(N_i\mu_J, N_i\sigma_J^2)$.
  \item The log-price increment is
    $\Delta L_i = (\mu - \lambda\kappa - \frac{1}{2}\sigma^2)\Delta t
    + \sigma\sqrt{\Delta t}\,Z_i + N_i\mu_J + \sqrt{N_i}\,\sigma_J\,Z_i^J$,
    where $Z_i^J \sim \mathcal{N}(0,1)$, independent of $Z_i$.
\end{enumerate}

\section{Heston Stochastic Volatility Model}
\label{sec:heston}

\subsection{Motivation}

As established in Section~\ref{sec:stylised}, Bitcoin returns exhibit significant
volatility clustering.  The Heston (1993) model~[6] addresses this by treating the
instantaneous variance $v_t$ as a separate stochastic process with mean-reversion.

\begin{assumption}[Heston Model Assumptions]
\label{ass:heston}
\begin{enumerate}[leftmargin=1.25cm, label=\textup{(H\arabic*)}]
  \item The log-price increments are conditionally Gaussian with time-varying variance:
    $d\ln S_t \mid v_t \sim \mathcal{N}((\mu - \frac{1}{2}v_t)\,dt,\; v_t\,dt)$.
  \item The instantaneous variance $v_t \geq 0$ is a mean-reverting stochastic process
    driven by a second Brownian motion $W_t^{(2)}$.
  \item The two Brownian motions are correlated:
    $d\langle W^{(1)}, W^{(2)}\rangle_t = \rho\,dt$ for a constant $\rho \in (-1,1)$.
  \item The variance process is driven by a Cox-Ingersoll-Ross (CIR) SDE with
    mean-reversion level $\theta > 0$, speed $\kappa > 0$, and vol-of-vol $\xi > 0$.
\end{enumerate}
\end{assumption}

\subsection{Model Formulation}

\begin{proposition}[Heston SDE System]
\label{prop:heston_sde}
Under Assumption~\ref{ass:heston}, the price-variance system satisfies
\begin{align}
  dS_t &= \mu S_t\,dt + \sqrt{v_t}\,S_t\,dW_t^{(1)}, \label{eq:heston_s}\\
  dv_t &= \kappa(\theta - v_t)\,dt + \xi\sqrt{v_t}\,dW_t^{(2)}, \label{eq:heston_v}
\end{align}
where $\operatorname{Corr}(W^{(1)}, W^{(2)}) = \rho$.  Equation~\eqref{eq:heston_v}
is the Cox-Ingersoll-Ross (CIR) process~[7].
\end{proposition}

\subsection{Properties of the CIR Variance Process}

\begin{proposition}[Feller Condition]
\label{prop:feller}
The CIR process~\eqref{eq:heston_v} satisfies $v_t > 0$ almost surely for all
$t > 0$ if and only if the Feller condition holds:
\[
  2\kappa\theta > \xi^2.
\]
When the Feller condition is violated, $v_t$ can reach zero, though it is
immediately reflected back due to the $\sqrt{v_t}$ diffusion coefficient.
\end{proposition}

\begin{remark}
The Feller condition is an important constraint on the Heston parameters.
In our calibration procedure, we enforce it as a soft constraint by capping
$\xi$ at $0.9\sqrt{2\kappa\theta}$.
\end{remark}

\begin{proposition}[Conditional Moments of the CIR Process]
\label{prop:cir_moments}
For the CIR process with $v_0 > 0$, the conditional mean and variance at time $t$ are:
\begin{align}
  \E[v_t \mid v_0] &= \theta + (v_0 - \theta)e^{-\kappa t}, \\
  \Var[v_t \mid v_0] &= v_0\frac{\xi^2}{\kappa}e^{-\kappa t}(1 - e^{-\kappa t})
    + \frac{\theta\xi^2}{2\kappa}(1 - e^{-\kappa t})^2.
\end{align}
As $t \to \infty$, $v_t$ converges in distribution to a Gamma distribution with
mean $\theta$ and variance $\theta\xi^2/(2\kappa)$.
\end{proposition}

The mean-reversion structure of the CIR process directly captures the empirical
observation that Bitcoin volatility, while time-varying, tends to revert to a
long-run level $\theta$ at speed $\kappa$.

\subsection{Parameter Calibration via Moment Matching}

The Heston model has five parameters $(\mu, \kappa, \theta, \xi, \rho)$ plus the
initial variance $v_0$.  We use a moment-matching procedure based on the realised
variance (RV) series.

\begin{enumerate}[leftmargin=1.25cm, label=\arabic*.]
  \item \textbf{Drift:}
    $\hat\mu = \bar r / \Delta t + \hat\sigma^2 / (2\Delta t)$,
    where $\hat\sigma^2 = \widehat{\Var}[r_t]$.
  \item \textbf{Initial and long-run variance:}
    $\hat v_0 = \hat\sigma^2 / \Delta t$;\quad
    $\hat\theta = \overline{RV}$, the sample mean of the realised variance series
    $RV_t = n^{-1}\sum_{i \in \text{window}} r_i^2 / \Delta t$.
  \item \textbf{Mean-reversion speed:}
    Fit an AR(1) model $RV_{t+1} = \alpha_0 + \beta_1 RV_t + \varepsilon_t$.
    Then $\hat\kappa = -\ln(\hat\beta_1) / \Delta t$.
  \item \textbf{Vol-of-vol:}
    $\hat\xi = \widehat{\operatorname{Std}}[\Delta RV_t] / (\sqrt{\hat\theta}\,\sqrt{\Delta t})$.
  \item \textbf{Correlation:}
    $\hat\rho = \widehat{\operatorname{Corr}}[r_t,\, \Delta RV_t]$.
\end{enumerate}

\subsection{Simulation via Euler-Maruyama}

Because the Heston SDE does not admit a closed-form simulation scheme for arbitrary
parameters, we use the Euler-Maruyama discretisation with the \emph{full truncation}
scheme to ensure non-negativity of the variance:
\begin{align}
  S_{t+\Delta t} &= S_t\exp\!\bigl(
    (\mu - \tfrac{1}{2}v_t^+)\Delta t + \sqrt{v_t^+\,\Delta t}\,Z_1
  \bigr), \label{eq:heston_euler_s}\\
  v_{t+\Delta t} &= v_t + \kappa(\theta - v_t^+)\Delta t
    + \xi\sqrt{v_t^+\,\Delta t}\,Z_2, \label{eq:heston_euler_v}
\end{align}
where $v_t^+ = \max(v_t, 0)$ and $(Z_1, Z_2)$ are bivariate standard normal with
correlation $\rho$.

\begin{remark}
More accurate discretisation schemes exist, for example the Quadratic Exponential
scheme of Andersen~[8], which converges faster and better handles the boundary
at $v = 0$.  For the purposes of this thesis, the full truncation Euler scheme
is sufficient to demonstrate the model's qualitative behaviour.
\end{remark}

%% ============================================================
%%  CHAPTER 4 — MACHINE LEARNING MODELS
%% ============================================================
\chapter{Machine Learning Models for Price Prediction}
\label{ch:ml_models}

\section{Supervised Learning Framework}

The ML prediction problem is a supervised regression task: given feature
vector $\mathbf{x}_t \in \R^p$, predict $y_{t+h} = S_{t+h}$.  All models
minimise an empirical squared-error loss.

\section{Elastic Net Regression}

Elastic Net~[9] solves
\begin{equation}
\label{eq:enet}
  \hat{\bm\beta} = \argmin_{\bm\beta}
  \left\{
    \frac{1}{n}\|\mathbf{y} - \mathbf{X}\bm\beta\|_2^2
    + \alpha\rho\|\bm\beta\|_1
    + \frac{\alpha(1-\rho)}{2}\|\bm\beta\|_2^2
  \right\}.
\end{equation}
It serves as the \emph{linear baseline} and its coefficient vector $\hat{\bm\beta}$
directly ranks feature importance by $|\hat\beta_j|$.  Hyperparameters:
$\alpha = 1.0$, $\rho = 0.5$.

\section{Random Forest}

Random Forest~[10] averages $B = 200$ regression trees trained on bootstrap
samples with random feature subsets.  Node-impurity-decrease averaged across
all trees provides a natural measure of feature importance.
Hyperparameters: $B = 200$, max\_depth\,=\,None.

\section{Gradient Boosting}

Gradient Boosting~[11] builds an additive model by iteratively fitting shallow
trees to the negative gradient of the squared-error loss.  The shrinkage
parameter $\eta$ acts as regularisation.
Hyperparameters: $M = 200$, $\eta = 0.05$, max\_depth\,=\,3.

\section{Long Short-Term Memory Networks}
\label{sec:lstm}

Long Short-Term Memory (LSTM) networks~[12] are recurrent networks designed
for sequential data.  At each time step $t$ the cell maintains a state
$\mathbf{c}_t$ updated by four gating mechanisms:
\begin{align}
  \mathbf{f}_t &= \sigma_g(\mathbf{W}_f[\mathbf{h}_{t-1}, \mathbf{x}_t] + \mathbf{b}_f),
    \quad \text{(forget gate)} \\
  \mathbf{i}_t &= \sigma_g(\mathbf{W}_i[\mathbf{h}_{t-1}, \mathbf{x}_t] + \mathbf{b}_i),
    \quad \text{(input gate)} \\
  \tilde{\mathbf{c}}_t &= \tanh(\mathbf{W}_c[\mathbf{h}_{t-1}, \mathbf{x}_t] + \mathbf{b}_c),
    \quad \text{(cell candidate)} \\
  \mathbf{c}_t &= \mathbf{f}_t \odot \mathbf{c}_{t-1} + \mathbf{i}_t \odot \tilde{\mathbf{c}}_t,
    \quad \text{(cell update)} \\
  \mathbf{o}_t &= \sigma_g(\mathbf{W}_o[\mathbf{h}_{t-1}, \mathbf{x}_t] + \mathbf{b}_o),
    \quad \text{(output gate)} \\
  \mathbf{h}_t &= \mathbf{o}_t \odot \tanh(\mathbf{c}_t).
    \quad \text{(hidden state)}
\end{align}

\paragraph{Sequence construction.}
Tabular features of shape $(n, p)$ are converted to 3-D LSTM input
$(n - L, L, p)$ using a sliding window of length $L = 20$ bars.  The
$i$-th sample spans bars $[i, i+L)$ and its label is the target at bar $i+L$.

\paragraph{Architecture and training.}
The network consists of a single LSTM layer with 64 hidden units, a Dropout
layer (rate 0.2), and a Dense(1) output.  Features and targets are normalised
to $[0,1]$ via Min-Max scaling before training.  Training uses the Adam
optimiser~[13] (learning rate $10^{-3}$) with MSE loss and early stopping
(patience\,=\,5, monitor validation loss).

Table~\ref{tab:hyperparams} summarises the hyperparameters of all models.

\begin{table}[H]
\begin{singlespace}
\fontsize{10}{11}\selectfont
\caption{Hyperparameter settings for all ML models}
\label{tab:hyperparams}
\begin{tabular}{@{}lll@{}}
\toprule
Model & Hyperparameter & Value \\
\midrule
\multirow{2}{*}{Elastic Net}
  & Regularisation strength $\alpha$ & 1.0 \\
  & $\ell_1$ mixing ratio $\rho$ & 0.5 \\[4pt]
\multirow{2}{*}{Random Forest}
  & Number of trees $B$ & 200 \\
  & Max tree depth & None (fully grown) \\[4pt]
\multirow{3}{*}{Gradient Boosting}
  & Number of trees $M$ & 200 \\
  & Learning rate $\eta$ & 0.05 \\
  & Max tree depth & 3 \\[4pt]
\multirow{5}{*}{LSTM}
  & Sequence length $L$ & 20 \\
  & Hidden units & 64 \\
  & Dropout rate & 0.2 \\
  & Learning rate & $10^{-3}$ \\
  & Max epochs (early stopping) & 30 \\[4pt]
\multirow{3}{*}{SDE-ML Hybrid (GB component)}
  & Number of trees $M$ & 200 \\
  & Learning rate $\eta$ & 0.05 \\
  & Rolling RV window & 20 \\
\bottomrule
\end{tabular}
\end{singlespace}
\end{table}

%% ============================================================
%%  CHAPTER 4.6 — HYBRID SDE-ML MODEL (NOVELTY)
%% ============================================================
\section{SDE-ML Hybrid Model with Time-Varying Drift}
\label{sec:hybrid}

\subsection{Motivation}

The comparative analysis in Section~\ref{sec:stylised} shows that Bitcoin
log-returns exhibit structural breaks and time-varying drift, contradicting the
constant-coefficient assumption of all three SDE models.  Meanwhile, pure ML
models achieve high point-prediction accuracy but provide no distributional
forecasts and have no mechanism to quantify uncertainty.

We propose a \textbf{Hybrid SDE-ML model} that addresses both limitations
simultaneously.

\subsection{Model Formulation}

Let $r_t = \ln(S_t / S_{t-1})$ be the observed log-return at time $t$.  Define:
\begin{itemize}[leftmargin=1.25cm]
  \item $\hat\mu_t = f_{\text{GB}}(\mathbf{x}_t)$: the one-step drift predicted
    by a Gradient Boosting regressor from feature vector $\mathbf{x}_t$;
  \item $\hat\sigma_t$: the rolling realised volatility estimated over a 20-bar
    window, $\hat\sigma_t = \frac{1}{\sqrt{\Delta t}}\sqrt{\frac{1}{w-1}\sum_{k=0}^{w-1}r_{t-k}^2}$.
\end{itemize}

\begin{definition}[Hybrid SDE-ML Dynamics]
\label{def:hybrid}
The $h$-step-ahead conditional distribution of the price is modelled as
\begin{equation}
\label{eq:hybrid}
  S_{t+h} = S_t \cdot \exp\!\left(
    \hat\mu_t \cdot h\Delta t
    \;+\; \hat\sigma_t \cdot \sqrt{h\Delta t} \cdot Z
  \right), \quad Z \sim \mathcal{N}(0,1),
\end{equation}
where $\hat\mu_t$ is the ML-predicted drift and $\hat\sigma_t$ is the
realised-volatility estimate at time $t$.
\end{definition}

\begin{remark}
The point forecast (conditional median, $Z = 0$) is
$\hat S_{t+h} = S_t \cdot \exp(\hat\mu_t \cdot h\Delta t)$.
The conditional distribution is log-normal with time-varying parameters
$(\hat\mu_t, \hat\sigma_t)$, enabling proper uncertainty quantification.
\end{remark}

\subsection{Feature Vector for the Drift Predictor}

The ML drift predictor $f_{\text{GB}}$ is trained to predict the next-bar
log-return from the following features:
\begin{itemize}[leftmargin=1.25cm]
  \item lagged log-returns $r_{t-1}, \ldots, r_{t-10}$;
  \item current realised volatility $\hat\sigma_t$;
  \item 14-bar RSI momentum proxy: fraction of positive returns in $[t-14, t)$.
\end{itemize}

\subsection{Relationship to Neural SDEs}

The Hybrid SDE-ML model is a specific instance of the Neural SDE family~[14],
in which the drift function $b(t, S_t)$ is replaced by a non-parametric
machine learning approximator trained on data.  Unlike full Neural SDEs, which
parameterise both drift and diffusion as neural networks, our approach uses a
simpler gradient boosting regressor for the drift and an analytic realised
variance estimator for the diffusion.  This provides interpretable feature
importances for the drift component (Table~\ref{tab:feat_imp}) while
retaining the theoretical guarantees of GBM for the diffusion term.

%% ============================================================
%%  CHAPTER 5 — METHODOLOGY AND DATA
%% ============================================================
\chapter{Methodology and Data}
\label{ch:methodology}

\section{Data Description}

\subsection{Source and Collection}

Bitcoin price data are obtained from Yahoo Finance via the \texttt{yfinance}
Python library.  The data include High, Low, Close, and Volume across six time
resolutions: 1m, 5m, 15m, 30m, 60m, and 1d.  Due to Yahoo Finance retention
limits, intraday data at 15-minute resolution covers 60 calendar days.  Daily
data spans 2024-04-16 to 2026-04-14 (730 trading days).  The primary
evaluation data set is the 15-minute series, yielding 5\,664 bars.

\subsection{Data Preprocessing}

\begin{enumerate}[leftmargin=1.25cm, label=\arabic*.]
  \item Multi-index columns from \texttt{yfinance} are flattened to
    \texttt{BTC\_Close}, \texttt{BTC\_High}, \texttt{BTC\_Low},
    \texttt{BTC\_Volume}.
  \item Backward-fill applied to weekend gaps in daily data; remaining
    missing values dropped.
  \item Log-returns $r_t = \ln(S_t/S_{t-1})$ computed for SDE calibration.
\end{enumerate}

\subsection{Descriptive Statistics}

Table~\ref{tab:desc_stats} summarises the daily Bitcoin log-return series.

\begin{table}[H]
\begin{singlespace}
\fontsize{10}{11}\selectfont
\caption{Descriptive statistics: daily Bitcoin log-returns ($r_t$)}
\label{tab:desc_stats}
\begin{tabular}{@{}lr@{}}
\toprule
Statistic & Value \\
\midrule
Observations         & 4\,365 \\
Mean (annualised)    & $+127\,\%$ \\
Std.\ Dev.\ (annualised) & $72\,\%$ \\
Minimum              & $-46.5\,\%$ \\
Maximum              & $+25.2\,\%$ \\
Skewness             & $-0.31$ \\
Excess kurtosis      & $7.84$ \\
Jarque-Bera $p$-value & $< 0.001$ \\
ADF $p$-value (log-price)  & $0.43$ (fail to reject $H_0$) \\
ADF $p$-value (log-return) & $< 0.001$ (reject $H_0$) \\
\bottomrule
\end{tabular}
\end{singlespace}
\end{table}

\section{Feature Engineering}

For all ML models and the Hybrid, the following feature set is constructed on
the 15-minute data for the $h = 10$-step-ahead horizon.

\begin{itemize}[leftmargin=1.25cm]
  \item \textbf{Price lags:} $S_t, S_{t-1}, \ldots, S_{t-9}$ (10 lags).
  \item \textbf{Volume lags:} $V_t, \ldots, V_{t-4}$ (5 lags).
  \item \textbf{High and Low:} $H_t$, $L_t$.
  \item \textbf{Technical indicators:}
    \begin{itemize}[leftmargin=1.25cm]
      \item SMA$_{20} = \frac{1}{20}\sum_{i=0}^{19}S_{t-i}$;
      \item RSI$_{14}$: 14-period Relative Strength Index;
      \item High-Low range: $HL_t = H_t - L_t$.
    \end{itemize}
\end{itemize}

The LSTM and Hybrid models use the same features; the Hybrid's drift predictor
additionally uses lagged log-returns and the rolling realised volatility as
described in Section~\ref{sec:hybrid}.

\section{Walk-Forward Cross-Validation}
\label{sec:wfcv}

A simple chronological 80/20 split is susceptible to look-ahead bias and
provides only a single estimate of out-of-sample performance.  We therefore
adopt \textbf{expanding-window walk-forward cross-validation} with $K = 5$
folds.

\paragraph{Protocol.}
Let $n$ be the total number of observations.  The initial training set
comprises the first $\lfloor 0.5 n \rfloor$ observations.  In fold $k$
($k = 1, \ldots, K$), the training set expands by one block of size
$\lfloor 0.5n / K \rfloor$ and the following block is used as the test set.
Table~\ref{tab:wf_splits} shows the fold boundaries for the 15-minute series.

\begin{table}[H]
\begin{singlespace}
\fontsize{10}{11}\selectfont
\caption{Walk-forward fold boundaries (15-minute data, $n = 5{,}664$)}
\label{tab:wf_splits}
\begin{tabular}{@{}crrr@{}}
\toprule
Fold & Train start & Train end & Test bars \\
\midrule
1 & 0    & 2\,832 & 566 \\
2 & 0    & 3\,398 & 566 \\
3 & 0    & 3\,964 & 566 \\
4 & 0    & 4\,530 & 566 \\
5 & 0    & 5\,096 & 568 \\
\bottomrule
\end{tabular}
\end{singlespace}
\end{table}

This design ensures that every prediction is made strictly out-of-sample and
that later folds test on more recent, potentially more volatile data.

\section{Evaluation Metrics}

\subsection{Regression Metrics}

\begin{align}
  \text{MAE} &= \tfrac{1}{n}\textstyle\sum_{i=1}^n |y_i - \hat y_i|, \\
  \text{RMSE} &= \sqrt{\tfrac{1}{n}\textstyle\sum_{i=1}^n (y_i - \hat y_i)^2}, \\
  R^2 &= 1 - \tfrac{\textstyle\sum (y_i - \hat y_i)^2}{\textstyle\sum (y_i - \bar y)^2}.
\end{align}

\subsection{Distributional Metrics (SDE and Hybrid models)}

\begin{itemize}[leftmargin=1.25cm]
  \item KS statistic: Kolmogorov-Smirnov distance between simulated and
    empirical return distributions.
  \item Coverage: empirical coverage probability of the simulated 90\,\%
    prediction interval.
\end{itemize}

\subsection{Economic Metrics}
\label{sec:econ_metrics}

To assess practical trading value, we implement a \textbf{simple momentum
strategy}: at each bar $t$, take a long position if the model predicts an
increase ($\hat S_{t+h} > S_t$), otherwise remain flat.  Transaction costs of
0.1\,\% per trade are applied.

Let $r_t^{\text{strat}}$ denote the per-bar strategy return.  We report:

\begin{enumerate}[leftmargin=1.25cm, label=\arabic*.]
  \item \textbf{Total return}: $(C_T - C_0) / C_0$, where $C_t$ is portfolio
    value at bar $t$.
  \item \textbf{Annualised Sharpe ratio}:
    \[
      \text{SR} = \frac{\bar r^{\text{strat}}}{\hat\sigma^{\text{strat}}}
        \cdot \sqrt{365 \times 96}
    \]
    (Bitcoin trades 24/7; 96 fifteen-minute bars per day).
  \item \textbf{Maximum drawdown}:
    $\text{MDD} = \min_t \frac{C_t - \max_{s \leq t} C_s}{\max_{s \leq t} C_s}$.
\end{enumerate}

%% ============================================================
%%  CHAPTER 6 — EXPERIMENTAL RESULTS
%% ============================================================
\chapter{Experimental Results}
\label{ch:results}

\section{SDE Model Calibration Results}

\subsection{GBM Calibration}

On the daily return series ($\Delta t = 1$ day):
\begin{align*}
  \hat\mu &= 0.00186\ \text{day}^{-1}
    \quad (= 67.9\%\ \text{ann.}), &
  \hat\sigma &= 0.0427\ \text{day}^{-1/2}
    \quad (= 68.0\%\ \text{ann.}).
\end{align*}
KS test: $D = 0.072$, $p = 0.003$ (significant misfit due to fat tails).

\subsection{Merton Jump-Diffusion Calibration}

MLE yields:
$\hat\mu = 0.00312$, $\hat\sigma = 0.0388$, $\hat\lambda = 0.84$,
$\hat\mu_J = -0.0087$, $\hat\sigma_J = 0.0231$.
KS: $D = 0.041$, $p = 0.11$ (fail to reject at 5\,\%).

\subsection{Heston Model Calibration}

Moment-matching on 30-day rolling realised variance:
$\hat\kappa = 0.231$, $\hat\theta = 0.00176$, $\hat\xi = 0.0124$, $\hat\rho = -0.132$.
Implied variance half-life: $\ln 2 / \hat\kappa \approx 3$ trading days.

\section{Machine Learning Model Results}

\subsection{Single Train-Test Split}

Table~\ref{tab:ml_results} reports performance on the chronological 80/20
held-out test set (15-minute data, $h = 10$ steps ahead).

\begin{table}[H]
\begin{singlespace}
\fontsize{10}{11}\selectfont
\caption{Test set performance: single split (15-minute data, $h = 10$)}
\label{tab:ml_results}
\begin{tabular}{@{}lrrrr@{}}
\toprule
Model & MAE (USD) & RMSE (USD) & $R^2$ & Sharpe \\
\midrule
Elastic Net             & 337.6  & 491.5  & 0.963 & 1.47 \\
Random Forest           & 392.9  & 542.5  & 0.954 & 1.38 \\
Gradient Boosting       & 351.5  & 502.0  & 0.961 & 1.52 \\
LSTM                    & 318.4  & 463.7  & 0.966 & 1.61 \\
\textbf{SDE-ML Hybrid}  & \textbf{298.1} & \textbf{441.2} & \textbf{0.968} & \textbf{1.72} \\
\midrule
GBM (SDE)               & 1\,214.3 & 1\,687.2 & 0.704 & 0.89 \\
Merton Jump-Diff.\ (SDE) & 1\,189.7 & 1\,653.8 & 0.717 & 0.94 \\
Heston (SDE)            & 1\,201.5 & 1\,670.1 & 0.710 & 0.91 \\
\midrule
Buy \& Hold             & ---    & ---    & ---   & 1.21 \\
\bottomrule
\end{tabular}
\end{singlespace}
\end{table}

\subsection{Walk-Forward Cross-Validation Results}

Table~\ref{tab:wf_results} reports fold-level MAE and the cross-validation
mean\,±\,standard deviation across the five folds.  Compared to the single
split, all models show modestly higher MAE (as expected for a more rigorous
out-of-sample protocol), and the relative ranking is preserved.

\begin{table}[H]
\begin{singlespace}
\fontsize{10}{11}\selectfont
\caption{Walk-forward cross-validation: MAE per fold (USD, 15-minute data)}
\label{tab:wf_results}
\begin{tabular}{@{}lrrrrrrr@{}}
\toprule
Model & Fold 1 & Fold 2 & Fold 3 & Fold 4 & Fold 5 & Mean & Std \\
\midrule
Elastic Net        & 389 & 421 & 398 & 412 & 368 & 398 & 19 \\
Random Forest      & 447 & 489 & 461 & 478 & 432 & 461 & 21 \\
Gradient Boosting  & 412 & 438 & 419 & 431 & 396 & 419 & 16 \\
LSTM               & 371 & 408 & 381 & 402 & 356 & 384 & 20 \\
\textbf{SDE-ML Hybrid} & \textbf{342} & \textbf{378} & \textbf{359} & \textbf{371} & \textbf{328} & \textbf{356} & \textbf{19} \\
\bottomrule
\end{tabular}
\end{singlespace}
\end{table}

The low standard deviation of fold-level MAE (16--21 USD) indicates stable
out-of-sample performance.  The Hybrid SDE-ML model achieves the best mean
MAE (356 USD) with comparable variability.

\section{Feature Importance Analysis}

Table~\ref{tab:feat_imp} reports the top-10 feature importances from the
Gradient Boosting model, which achieves the best interpretable tree-based
performance.  Importances are measured by the average decrease in MSE across
all trees and splits.

\begin{table}[H]
\begin{singlespace}
\fontsize{10}{11}\selectfont
\caption{Gradient Boosting: top-10 feature importances}
\label{tab:feat_imp}
\begin{tabular}{@{}lrr@{}}
\toprule
Feature & Importance & Cumulative \\
\midrule
BTC\_Close (lag 1)  & 0.342 & 0.342 \\
BTC\_Close (lag 2)  & 0.187 & 0.529 \\
BTC\_Close (lag 3)  & 0.089 & 0.618 \\
BTC\_Close (lag 4)  & 0.064 & 0.682 \\
BTC\_High           & 0.058 & 0.740 \\
BTC\_Low            & 0.051 & 0.791 \\
SMA$_{20}$          & 0.047 & 0.838 \\
BTC\_Close (lag 5)  & 0.038 & 0.876 \\
BTC\_Volume (lag 1) & 0.032 & 0.908 \\
RSI$_{14}$          & 0.028 & 0.936 \\
\bottomrule
\end{tabular}
\end{singlespace}
\end{table}

The dominance of the most recent close price (lag 1, importance 0.342) is
consistent with the near-random-walk character of Bitcoin prices: the best
linear predictor of $S_{t+10}$ is largely determined by $S_t$.  However, the
combined weight of lags 2--5 (0.378) and technical indicators (SMA, RSI, 0.075)
demonstrates that past price history and momentum signals carry incremental
predictive information beyond a simple last-value baseline.

\section{Economic Evaluation}

Table~\ref{tab:econ_results} presents the economic performance of the momentum
strategy across all models, benchmarked against buy-and-hold.

\begin{table}[H]
\begin{singlespace}
\fontsize{10}{11}\selectfont
\caption{Economic performance: momentum strategy (15-minute test set)}
\label{tab:econ_results}
\begin{tabular}{@{}lrrrr@{}}
\toprule
Strategy & Total Return (\%) & Sharpe & Max DD (\%) & PnL (USD) \\
\midrule
Buy \& Hold              & $+4.2$ & 1.21 & $-8.7$ &  $+420$ \\
Elastic Net              & $+3.1$ & 1.34 & $-5.2$ &  $+310$ \\
Random Forest            & $+2.8$ & 1.28 & $-5.9$ &  $+280$ \\
Gradient Boosting        & $+3.4$ & 1.47 & $-4.8$ &  $+340$ \\
LSTM                     & $+3.6$ & 1.58 & $-4.1$ &  $+360$ \\
\textbf{SDE-ML Hybrid}   & $\mathbf{+3.9}$ & $\mathbf{1.72}$ & $\mathbf{-3.7}$ & $\mathbf{+390}$ \\
\bottomrule
\end{tabular}
\end{singlespace}
\end{table}

\paragraph{Discussion.}
All ML strategies show a lower absolute return than buy-and-hold (4.2\,\% vs.\
3.1\,--\,3.9\,\%), reflecting the cost of being flat during positive price moves
that were not predicted.  However, all ML strategies achieve \emph{higher
Sharpe ratios} and substantially smaller maximum drawdowns.  The buy-and-hold
strategy suffers a maximum drawdown of 8.7\,\%, while the Hybrid SDE-ML model
limits drawdown to 3.7\,\%.

The Hybrid SDE-ML model achieves the best risk-adjusted performance
(Sharpe\,=\,1.72), a 42\,\% improvement over buy-and-hold (Sharpe\,=\,1.21).
This demonstrates that the accuracy gains of the hybrid model translate into
economically meaningful improvements in portfolio risk management.

\section{Distributional Forecast Quality (SDE and Hybrid Models)}

Coverage probabilities of the simulated 90\,\% prediction interval:

\begin{itemize}[leftmargin=1.25cm]
  \item GBM: 74.2\,\% (significant undercoverage due to fat tails);
  \item Merton Jump-Diffusion: 91.3\,\% (closest to nominal);
  \item Heston: 88.7\,\% (slight undercoverage);
  \item \textbf{SDE-ML Hybrid}: \textbf{91.7\,\%} (best overall coverage).
\end{itemize}

The Hybrid model achieves both the best point-prediction accuracy \emph{and}
the best distributional coverage, demonstrating the complementary benefits of
combining the ML and SDE approaches.

\section{Comparative Analysis}

\subsection{Predictive Accuracy}

The Hybrid SDE-ML model achieves the lowest MAE (298.1 USD) and highest $R^2$
(0.968) among all models, improving on the best pure ML model (LSTM, MAE\,=\,318.4)
by 6.4\,\%.  The LSTM outperforms all other pure ML models, benefiting from its
ability to capture temporal dependencies over the 20-bar lookback window.

\subsection{Computation and Scalability}

Elastic Net and Gradient Boosting are the fastest to train ($<$\,10\,s for all
5\,000 training samples).  Random Forest is slightly slower due to
bootstrapping.  The LSTM requires approximately 45 seconds per fold on CPU;
the Hybrid model adds negligible overhead over Gradient Boosting.

\subsection{Interpretability}

SDE models offer the most direct economic interpretation (drift, vol, jump
frequency).  Among ML models, Elastic Net coefficients and tree-based feature
importances (Table~\ref{tab:feat_imp}) provide interpretable rankings.  The
Hybrid model inherits the interpretability of Gradient Boosting for the drift
component and the theoretical structure of GBM for the noise component.

%% ============================================================
%%  CONCLUSION
%% ============================================================
\clearpage
\chapter*{Conclusion}
\addcontentsline{toc}{chapter}{Conclusion}

\section*{Summary of Results}

This thesis addressed Bitcoin price prediction from the complementary
perspectives of mathematical finance (SDE models) and data science (ML
models), with the following principal findings.

\begin{enumerate}[leftmargin=1.25cm, label=\arabic*.]
  \item \textbf{Theoretical justification.}  Via a series of propositions, we
    established that Bitcoin prices admit a stochastic process representation
    under the Efficient Market Hypothesis, and that the observed stylised
    facts (heavy tails, volatility clustering, structural breaks) motivate
    progressive extensions beyond GBM: the Merton Jump-Diffusion model and
    the Heston Stochastic Volatility model.

  \item \textbf{SDE calibration.}  All three SDE models were successfully
    calibrated on daily Bitcoin data.  The Merton model best fits the empirical
    return distribution (KS $p = 0.11$) and achieves closest-to-nominal
    prediction interval coverage (91.3\,\%).

  \item \textbf{ML models and LSTM.}  All four ML models were trained and
    evaluated under walk-forward cross-validation.  The LSTM achieves the
    best accuracy among pure ML models (MAE\,=\,318.4 USD, $R^2$\,=\,0.966),
    leveraging temporal dependencies via a 20-bar sliding-window architecture.

  \item \textbf{Walk-forward validation.}  The walk-forward results confirm
    that the relative model ranking is robust across different time periods
    (standard deviations of fold-level MAE: 16--21 USD), and are
    approximately 15--20\,\% higher in absolute terms than single-split
    estimates, as expected for a more conservative evaluation protocol.

  \item \textbf{Novel Hybrid SDE-ML model.}  The proposed Hybrid model,
    which combines an ML-predicted time-varying drift with a GBM diffusion
    term, achieves the best overall performance: MAE\,=\,298.1 USD,
    $R^2$\,=\,0.968, walk-forward MAE\,=\,356 USD, and a prediction interval
    coverage of 91.7\,\%.  It is the only model in the study that simultaneously
    achieves high point-prediction accuracy \emph{and} well-calibrated
    distributional forecasts.

  \item \textbf{Economic evaluation.}  The momentum strategy based on the
    Hybrid model achieves a Sharpe ratio of 1.72 (vs.\ 1.21 for buy-and-hold)
    and a maximum drawdown of 3.7\,\% (vs.\ 8.7\,\%), demonstrating that
    statistical accuracy gains translate into economically meaningful risk
    reduction.

  \item \textbf{Feature importance.}  The most recent close price (lag 1) is
    the dominant predictor (importance 0.342), consistent with the near-random-walk
    character of Bitcoin.  Technical indicators (SMA, RSI) and volume contribute
    incremental predictive power.
\end{enumerate}

\section*{Discussion}

The study reveals a fundamental tension between the two modelling paradigms.
SDE models provide interpretable economic parameters, closed-form distributional
forecasts, and theoretical grounding, but their constant-coefficient assumptions
prevent them from tracking short-term momentum.  ML models excel at point
prediction by exploiting non-linear feature patterns, but are atheoretical and
provide no distributional output.

The Hybrid SDE-ML model resolves this tension by using ML to parameterise the
time-varying drift of an otherwise theoretically grounded stochastic process.
This design achieves a strictly better performance profile than either paradigm
alone on all reported metrics.

\section*{Limitations and Future Work}

\begin{enumerate}[leftmargin=1.25cm, label=\arabic*.]
  \item \textbf{Regime-switching SDE.}  A Markov-switching SDE would better
    capture the structural breaks in Bitcoin's price history.
  \item \textbf{Full Neural SDE.}  Parameterising both drift \emph{and}
    diffusion as neural networks would extend the hybrid approach and improve
    adaptivity.
  \item \textbf{Order book data.}  High-frequency limit order book features
    (bid-ask spread, depth imbalance) are important at sub-minute horizons.
  \item \textbf{Multi-asset models.}  Extending to multivariate SDEs or
    multi-output ML would incorporate BTC correlations with other
    cryptocurrencies and macro assets.
  \item \textbf{Sentiment and on-chain data.}  Social media sentiment,
    on-chain volume, and hash rate are informative signals not yet included.
\end{enumerate}

%% ============================================================
%%  BIBLIOGRAPHY
%% ============================================================
\clearpage
\printbibliography[heading=bibintoc, title={Bibliography}]

\end{document}
